% ATTEMPT_STATUS: unresolved
% TOP_PROBLEM: {"problemId": "problem.bunyakovsky-conjecture", "canonicalTitle": "Bunyakovsky conjecture", "releaseRank": 102, "primaryDomain": "number_theory_arithmetic_geometry", "primaryDomainLabel": "Number theory & arithmetic geometry", "displayStatus": "open_disputed_claim", "releaseStatus": "open_disputed_claim"}
% !TeX program = lualatex
% Complete problem section copied verbatim from research_batch_101_103.tex.
\documentclass[11pt]{article}
\usepackage[margin=25mm]{geometry}
\usepackage{fontspec}
\setmainfont{Latin Modern Roman}
\usepackage{amsmath,amssymb,mathtools}
\usepackage{unicode-math}
\setmathfont{Latin Modern Math}
\usepackage{enumitem,needspace}
\usepackage{xurl,hyperref}
\hypersetup{colorlinks=true,urlcolor=blue}
\setcounter{secnumdepth}{1}
\setlength{\parindent}{0pt}
\setlength{\parskip}{5pt}
\setlength{\emergencystretch}{3em}
\setlist[itemize]{leftmargin=3em,itemsep=5pt,topsep=4pt,parsep=0pt}
\allowdisplaybreaks
\newcommand{\N}{\mathbb{N}}
\newcommand{\Z}{\mathbb{Z}}
\newcommand{\Q}{\mathbb{Q}}
\newcommand{\R}{\mathbb{R}}
\newcommand{\C}{\mathbb{C}}
\newcommand{\F}{\mathbb{F}}
\newcommand{\A}{\mathbb{A}}
\newcommand{\PP}{\mathbb P}
\newcommand{\E}{\mathbb{E}}
\newcommand{\eps}{\varepsilon}
\newcommand{\dd}{\,\mathrm d}
\newcommand{\sref}[2]{\hyperlink{source:#1:#2}{[S#2]}}
\newcommand{\eref}[2]{\hyperlink{extra:#1:#2}{[E#2]}}
% Turn the next switch off for a shorter reading copy. The quotations preserve
% every exactTarget field, including qualifications omitted from a short summary.
\newif\ifshowcatalogue
\showcataloguetrue
\newcommand{\cataloguescope}[1]{\ifshowcatalogue
 \paragraph{Exact scope in the supplied catalogue.}
 \begin{quote}\small #1\end{quote}\fi}
\begin{document}
\setcounter{section}{101}
% ================================================================
% problemId: problem.bunyakovsky-conjecture
% source releaseRank: 102; source releaseStatus: open_disputed_claim
% exactTarget SHA-256: 62e270d269e82f53d0bcc187ac56a1e341785d173040d79276b391c1c38a7b76
\clearpage
\section{\texorpdfstring{Bunyakovsky conjecture}{Bunyakovsky conjecture}}\label{102}
\subsection{Definitions and mathematical statement}
For a nonconstant integer polynomial $f$, a fixed prime divisor is a prime $p$ dividing $f(n)$ for every integer $n$. Suppose $f$ is irreducible over $\Q$, primitive, and has positive leading coefficient, with
\[
 \forall p\in\mathcal P\ \exists n\in\Z:\ p\nmid f(n).
\]
The assertion is that the set $\{n\in\Z_{>0}:f(n)\in\mathcal P\}$ is infinite. Positive leading coefficient ensures that signs cannot obstruct eventual positive prime values. No explicit count or density is required.
\par\smallskip\textit{Formulation and concept references:} \sref{102}{1}.
\cataloguescope{Prove or disprove that every irreducible polynomial f in Z[x] with positive leading coefficient and no fixed prime divisor takes prime values for infinitely many positive integer inputs.}
\subsection{Short English statement}
Must every irreducible positive-leading integer polynomial without a fixed prime divisor produce infinitely many primes?
% BEGIN RESEARCH ADDITION ATTEMPT-102
\subsection{Research attempt}

\paragraph{Current frontier.}
The degree-one case follows from Dirichlet's theorem. For an admissible
irreducible quadratic, Lemke Oliver's Theorem 1 proves infinitely many
squarefree values with at most two prime factors, extending Iwaniec's
$n^2+1$ result to general quadratics. ``At most two'' cannot be replaced by
``one'' in that theorem. The linear-sieve lower-bound function appearing in
his Lemma 2 is identically zero for sieve parameter $0<s\le2$.
\eref{102}{1}

The May 2026 revision of Browning--Sofos--Ter\"av\"ainen establishes
Bateman--Horn and polynomial Chowla results for almost all polynomials in
specified coefficient boxes, with quantitative exceptional sets
(Theorems 1.2 and 1.5). This is not an assertion for each fixed admissible
polynomial. \eref{102}{2}
A different recent advance, by Grimmelt and Merikoski, reaches exponent
$1.312$ for the greatest prime factor of $n^2+1$ infinitely often. A large
prime factor need not be the whole value. \eref{102}{3}
These results motivate combining roughness with a parity-sensitive statistic,
but none supplies the individual estimate required below.

\paragraph{Attack route.}
The local route would try to avoid successively more small prime divisors
using the Chinese remainder theorem; it needs uniform control as the modulus
grows with the input range. The coefficient-averaging route needs to transfer
an almost-everywhere theorem to a specified polynomial, including all its
translates. The selected route first removes values with three or more prime
factors, then seeks a strictly positive fraction of odd-parity values among
what remains. This isolates two separate missing estimates rather than hiding
the parity problem inside a heuristic independence assertion.

\paragraph{Attempt: exact local information.}
Fix one polynomial $f$ satisfying the original hypotheses, of degree $d$.
For an integer $q\ge1$, write
\[
 \rho_f(q)=\#\{a\bmod q:f(a)\equiv0\pmod q\}.
\]
For squarefree $q$ the Chinese remainder theorem gives
$\rho_f(q)=\prod_{p\mid q}\rho_f(p)$, and admissibility says
$\rho_f(p)<p$. If $X$ is a positive integer, residue-class counting gives
\begin{equation}\label{ra102:local}
 A_q(X):=\#\{X<n\le2X:q\mid f(n)\}
       =\frac{X\rho_f(q)}q+O(\rho_f(q)),
\end{equation}
with an absolute implied constant. This holds for every $q$, not just small
ones: each of its $\rho_f(q)$ residue classes contains $X/q+O(1)$ integers
in the interval.

For a \emph{fixed} squarefree $M$, the proportion of integers avoiding all
prime divisors of $M$ in their $f$-value is therefore
\[
 \prod_{p\mid M}\left(1-\frac{\rho_f(p)}p\right)>0.
\]
This positive quantity depends on $M$. Neither its positivity nor
\eqref{ra102:local} licenses sending $M\to\infty$ before controlling the
error and the shrinking density. Thus the simplest proposed local-to-global
proof already has an identified nonuniform limit.

\paragraph{An exact prime detector on rough values.}
Take $X$ large enough that $f(n)>1$ throughout $X<n\le2X$, and put
\[
 M_X=\max_{X<n\le2X}f(n),\qquad z_X=2M_X^{1/3},\qquad
 P(z)=\prod_{p<z}p.
\]
Since $f$ is fixed and has positive leading coefficient,
$M_X\asymp_f X^d$ and $\min_{X<n\le2X}f(n)\asymp_f X^d$.
In particular, after increasing $X$,
\begin{equation}\label{ra102:range}
 1<z_X<\min_{X<n\le2X}f(n),\qquad z_X^3=8M_X>M_X.
\end{equation}
Define the rough set and two statistics by
\begin{align*}
 \mathcal R_f(X)&=\{n\in\Z:X<n\le2X,\ (f(n),P(z_X))=1\},\\
 A_f(X)&=|\mathcal R_f(X)|,\qquad
 L_f(X)=\sum_{n\in\mathcal R_f(X)}\lambda(f(n)),
\end{align*}
where $\lambda(m)=(-1)^{\Omega(m)}$ is the Liouville function and
$\Omega$ counts prime factors \emph{with multiplicity}.

\textbf{Rough-parity lemma.} With these definitions, for every sufficiently
large integer $X$,
\begin{equation}\label{ra102:identity}
 \#\{X<n\le2X:f(n)\text{ is prime}\}
       =\frac{A_f(X)-L_f(X)}2.
\end{equation}
\emph{Proof.} Every prime factor of a rough value is at least $z_X$.
If it had three prime factors counted with multiplicity, its value would be
at least $z_X^3>M_X$, a contradiction. Thus every rough value has
$\Omega=1$ or $2$; the value $1$ has been excluded explicitly.
On these two possibilities, $(1-\lambda)/2$ is respectively $1$ or $0$.
Conversely, every prime value in the interval exceeds $z_X$ by
\eqref{ra102:range}, so none is lost by restricting to $\mathcal R_f(X)$.
Summing proves the identity. Prime squares have $\Omega=2$ and contribute
zero, as required.\hfill$\square$

A concrete sufficient pair of estimates is the following: for this fixed
$f$, there exist $c_f>0$ and $\delta_f>0$ such that, on an unbounded sequence
of dyadic values $X=2^j$,
\begin{equation}\label{ra102:missing}
 A_f(X)\ge c_f\frac{X}{\log X},\qquad
 L_f(X)\le(1-\delta_f)A_f(X).
\end{equation}
Equation \eqref{ra102:identity} then gives at least
$\delta_fc_fX/(2\log X)$ prime-producing inputs in each such block.
This would prove the required infinitude for $f$; establishing the estimates
for every admissible $f$ would prove the full target.

The estimates in \eqref{ra102:missing} are \emph{sufficient, not asserted
equivalent} to the original weak infinitude statement. The exact weak
condition is only $A_f(X)-L_f(X)>0$ in arbitrarily large dyadic blocks.
A positive density on the scale $X/\log X$ is an additional ambition of this
route, not a silent replacement of the problem.

\paragraph{Testing the rough-count estimate against available local data.}
Because $f$ is primitive, its reduction modulo a prime is not the zero
polynomial. Thus $\rho_f(p)\le d$ for every prime, including primes where
the degree drops. For squarefree $q$,
$\rho_f(q)\le d^{\omega(q)}=\tau_d(q)$, where $\tau_d$ counts ordered
factorizations into $d$ positive integers. Consequently
\begin{equation}\label{ra102:remainder}
 \sum_{\substack{q\le D\\q\text{ squarefree}}}
 \left|A_q(X)-\frac{X\rho_f(q)}q\right|
 \ll_d D(1+\log D)^{d-1}.
\end{equation}
Indeed, the sum of $\tau_d(q)$ up to $D$ counts tuples with product at most
$D$. Summing the last coordinate first bounds that count by
$D\bigl(\sum_{m\le D}m^{-1}\bigr)^{d-1}$.
Hence the elementary residue-class information gives a useful level
$D=X/(\log X)^B$ if $B$ is sufficiently large, for instance $B>d+1$.

At our cutoff $z_X\asymp_f X^{d/3}$, however,
\[
 \frac{\log D}{\log z_X}\longrightarrow\frac3d.
\]
Already in degree two this tends to $3/2$. Even granting the appropriate
one-dimensional sieve hypotheses, the linear lower sieve has zero main
lower-bound function there: its $f_{\mathrm{sieve}}(s)$ vanishes throughout
$0<s\le2$. This precise obstruction is visible in \eref{102}{1}, Lemma 2.
For $d\ge3$ the elementary level is no better relative to this cutoff.
Equation \eqref{ra102:remainder} is therefore not a proof of the first
estimate in \eqref{ra102:missing}.

The known quadratic almost-prime theorem uses substantially more structure
than this argument. Moreover, an arbitrary $P_2$ value may have a small
prime factor and need not belong to our particular rough set. Thus even an
infinite supply of $P_2$ values cannot simply be inserted as a lower bound
for $A_f(X)$. The squarefree convention in the cited theorem and the
multiplicity convention in our detector are also deliberately distinguished.
\eref{102}{1}

\paragraph{Where parity enters after exact inclusion--exclusion.}
The second estimate in \eqref{ra102:missing} asks for cancellation on a
thin, arithmetically selected subset. Expanding its roughness condition gives
an exact finite identity:
\begin{equation}\label{ra102:mobius}
 L_f(X)=\sum_{q\mid P(z_X)}\mu(q)
           \sum_{\substack{X<n\le2X\\q\mid f(n)}}\lambda(f(n)).
\end{equation}
For a root $a\bmod q$, choose its representative $0\le a<q$ and define
\[
 F_{q,a}(t)=\frac{f(a+qt)}q\in\Z[t].
\]
Integrality follows because the constant coefficient $f(a)$ is divisible by
$q$, and every positive-degree coefficient in $f(a+qt)$ is also divisible
by $q$. On the integers $m$ for which $X<a+qm\le2X$,
$f(a+qm)=qF_{q,a}(m)>0$. Complete multiplicativity gives
$\lambda(f(a+qm))=\lambda(q)\lambda(F_{q,a}(m))$.
All $q\mid P(z_X)$ are squarefree, so $\mu(q)\lambda(q)=1$.
Thus \eqref{ra102:mobius} becomes
\begin{equation}\label{ra102:quotient}
 L_f(X)=
 \sum_{q\mid P(z_X)}
 \ \sum_{\substack{0\le a<q\\ f(a)\equiv0\pmod q}}
 \ \sum_{\substack{m\in\Z\\X<a+qm\le2X}}
        \lambda(F_{q,a}(m)).
\end{equation}
There is no exchange of infinite sums: $P(z_X)$ is finite for each $X$.
The cancellation expected from the M\"obius signs has disappeared on pulling
out the factor $q$. The burden has moved to Liouville cancellation for a
large collection of quotient polynomials, often on very short intervals.
This identity does not claim that these new terms have the same sign; it
shows precisely where the signs still need to be controlled.

Each $F_{q,a}$ remains irreducible over $\Q$, since a nonconstant affine
substitution and division by a nonzero rational number preserve
irreducibility. Its coefficients and interval length nevertheless depend on
$q,a,X$, and it need not remain primitive or have no fixed prime divisor.
A fixed-polynomial asymptotic with an unspecified onset is not uniform
information for \eqref{ra102:quotient}. For moduli larger than the input
interval, even the inner intervals can contain at most one integer.

This also identifies why a hypothetical unweighted estimate
$\sum_{X<n\le2X}\lambda(f(n))=o(X)$ would not, by itself, settle the chosen
route. When $A_f(X)$ is only of size $X/\log X$, an error $o(X)$ can be larger
than the entire rough set. As a logical stress test, an abstract sequence of
signs can be $+1$ on any specified $o(X)$-element set and balanced on the
complement, while having total sum zero. This is not a constructed Liouville
counterexample; it refutes the inference from unweighted mean cancellation
to cancellation on every sparse selected set.

\paragraph{A counterexample search that fails for an explicit reason.}
One way to try to disprove the conjecture is to prescribe a different
obstructing prime at every small input. This succeeds for arbitrarily long
finite prefixes, but not with a single fixed polynomial for all inputs.
Here is the exact construction. Given $L\ge1$, choose distinct odd primes
$p_1,\ldots,p_L>L$ and use the Chinese remainder theorem to choose $C>0$
with
\[
 C\equiv-j^2\pmod{p_j}\quad(1\le j\le L),\qquad
 C>\max_j p_j^2.
\]
Then $f_C(t)=t^2+C$ is monic, irreducible over $\Q$ because it has no real
root, and has no fixed prime divisor. For $p>2$ a nonzero quadratic has at
most two roots modulo $p$; for $p=2$, the values at $0$ and $1$ have opposite
parities. Nevertheless $p_j\mid f_C(j)$ and $f_C(j)>p_j$ for every
$1\le j\le L$, so all these values are composite.

For example, exact CRT computation gives $C=29043056844$ for $L=8$ with
\[
 (p_1,\ldots,p_8)=(11,13,17,19,23,29,31,37).
\]
The accompanying script records each integer quotient $f_C(j)/p_j$ and
checks divisibility. The construction proves that there is no bound depending
only on the degree for the first positive prime-producing input. It does
\emph{not} contradict Bunyakovsky: the coefficient $C$ changes with $L$.
There is no compactness principle that turns these unbounded integer
coefficients into one fixed integer polynomial with all values composite.

\paragraph{Stress test: quantifiers and coefficient averages.}
A useful elementary equivalence makes the averaging issue particularly clear.
Suppose one could prove that \emph{every} admissible polynomial has at least
one prime value at a positive integer. For any fixed $f$ and any positive
integer $m$, the translate $f_m(t)=f(t+m)$ has the same irreducibility,
positive leading coefficient, and local admissibility. A prime value of
$f_m$ supplies a prime value of $f$ at an input greater than $m$. The
universal one-value assertion would therefore already imply infinitely many
values for every $f$.

The missing word is ``every.'' A polynomial and all its translates may lie
in the exceptional sets of a coefficient-box theorem. Quantitative bounds
for a high-dimensional box do not prevent a much thinner translation family
from being entirely exceptional. The hypotheses and exceptional-set conclusion
of \eref{102}{2}, Theorem 1.2, do not provide this transfer. No almost-sure
statement is being promoted to a universal one here.

\paragraph{Finite computations.}
An exact factor sieve was run on $n^2+1$ for $1\le n\le20000$. Every possible
prime factor up to $20000$ was processed through roots of $n^2\equiv-1$
modulo that prime; any remaining factor is prime. Independent factorization
checks were made on the deterministic sample specified in the script.
With the exact cutoff test $\ell^3\ge8M_X$, where $\ell$ is the least prime
factor, the following values were obtained:
\begin{center}
\begin{tabular}{r|r r|r r}
$X$ & $A_f(X)$ & $L_f(X)$ & prime values & rough $\Omega=2$ values\\
\hline
$100$   & $16$   & $-14$  & $15$  & $1$\\
$1000$  & $132$  & $-62$  & $97$  & $35$\\
$10000$ & $1047$ & $-389$ & $718$ & $329$
\end{tabular}
\end{center}
These are input counts on $X<n\le2X$, not counts up to $X$.
They check \eqref{ra102:identity}, its endpoints, and its multiplicity
convention on finite examples. They give no uniform estimate, no proof that
these ratios persist, and no information about all admissible polynomials.
The script \texttt{research\_batch\_101\_103\_checks.py} reproduces the table
without floating-point root comparisons.

\paragraph{Outcome.}
\textbf{Outcome: Conditional reduction.}

The exact rough-parity detector \eqref{ra102:identity}, the quotient-polynomial
identity \eqref{ra102:quotient}, and the arbitrary composite-prefix construction
have been proved. The first two expose a precise sufficient route through
\eqref{ra102:missing}; the third defeats a naive finite-search or
coefficient-compactness argument. These elementary deductions are not claimed
to be new results in the literature.

Neither the required rough-count lower bound nor the needed parity estimate
has been established for every fixed admissible $f$. In particular, the
computation does not prove infinitude even for $f(t)=t^2+1$. The approach
would need a genuinely parity-sensitive individual estimate on the rough
set, with errors small compared with that set and with the stated quantifiers.
No counterexample to the conjecture was found.

% END RESEARCH ADDITION ATTEMPT-102
\subsection{Sources}
\begin{itemize}\raggedright
\item[\textnormal{[S1]}]\hypertarget{source:102:1}{} Eric W. Weisstein with Ed Pegg Jr., Bouniakowsky Conjecture, MathWorld--A Wolfram Resource.
\url{https://mathworld.wolfram.com/BouniakowskyConjecture.html}.
{\small\textit{Catalogue formulation source.}}
% BEGIN RESEARCH ADDITION REFERENCES-102
\item[\textnormal{[E1]}]\hypertarget{extra:102:1}{}
R. J. Lemke Oliver, \emph{Almost-primes represented by quadratic polynomials},
Acta Arithmetica \textbf{151} (2012), no.~3, 241--261.
Theorem 1 (squarefree $P_2$ values for admissible quadratics), and Section 2.2,
Lemma 2 (the linear-sieve functions and bilinear remainder).
\url{https://lemkeoliver.github.io/papers/04-Quadratic.pdf}.
{\small\textit{Author's manuscript consulted 25 September 2026.}}
\item[\textnormal{[E2]}]\hypertarget{extra:102:2}{}
T. Browning, E. Sofos, and J. Ter\"av\"ainen,
\emph{Bateman--Horn, polynomial Chowla and the Hasse principle with
probability 1}, arXiv:2212.10373v2, 21 May 2026.
Theorems 1.2 and 1.5, including the coefficient-box and exceptional-set
quantifiers. \url{https://arxiv.org/html/2212.10373v2}.
{\small\textit{Consulted 25 September 2026; results averaged over
polynomials are distinguished from an assertion for every fixed polynomial.}}
\item[\textnormal{[E3]}]\hypertarget{extra:102:3}{}
L. Grimmelt and J. Merikoski,
\emph{On the greatest prime factor and uniform equidistribution of quadratic
polynomials}, arXiv:2505.00493 (2025), abstract and introduction,
the $n^{1.312}$ greatest-prime-factor bound for $n^2+1$.
\url{https://arxiv.org/abs/2505.00493}.
{\small\textit{Consulted 25 September 2026; a greatest-prime-factor result,
not a prime-values theorem.}}
% END RESEARCH ADDITION REFERENCES-102
\end{itemize}

\end{document}
